% No preamble. Prefix lpd-. Conditional downstream theorem, 17 October 2026.
% The actual uniform coefficient comparison is a named hypothesis here.
\section{From a uniform positive coefficient to finite local trajectories}
\label{lpd-section}

This section isolates the exact downstream consequence of a uniform
actual determinant comparison. The finite compactness, moment passage,
activity induction and endpoint shooting inputs are those of
Notebooks 31--32. The coefficient comparison below is an explicit
hypothesis of this section; it must be supplied by the actual scalar
and vector comparison, with the stated uniformity. No perturbative
reference coefficient is substituted for an unidentified actual one.

\subsection{Fixed data and the uniform coefficient hypothesis}
\label{lpd-data}

Fix the elementary quantum ratio \(R>1\), grouping \(J\), field
radius \(b>0\), carrier exponent \(a_0\), and the real extracted
finite activity system. Let \(q_a\) be its amplitude block length,
\(r_b=K_{\mathrm{amp}}b^2\) its endpoint ball radius, and
\(C_{\mathrm{int}}b^2\) its intermediate bound. These data are fixed
before the longer block length used below. The activity induction
holds for arbitrary prescribed couplings in one fixed weak range,
provided the retained periods exceed its threshold.

All finite extraction functionals \(\beta_j\) are the actual
short-lift functionals normalized by the actual classical reference.
For the actual finite local map \(F_{j,g}\), write
\begin{equation}
 d_j(g,x)=\beta_{j+1}(F_{j,g}(x)),\qquad
 x_{j+1}=P_{j+1}F_{j,g_j}(x_j),\qquad
 \gamma_{j+1}=\gamma_j-d_j(g_j,x_j).
 \label{lpd-updates}
\end{equation}
Real initial classes and the real branches are understood throughout.
The initial class may be zero, or any fixed real extracted class in
the endpoint ball. The actual coupling relation is
\(g_j=\gamma_j^{-1/2}\) whenever \(\gamma_j>0\).

Let \(D_{M,K}^{*}\) denote the actual leading determinant correction
at preceding homogeneous depth \(M\) and endpoint ratio \(K\) in a
joint infinite structural system. A star labels that structural
system, not a preferred limit of one fixed old depth. Its Maxwell
functional is the true infinite one, with the same classical
normalization used in the finite moment passage. The system includes
the shifted depths and endpoint ratios needed by the determinant
cocycle. In particular all comparisons below refer to the same
joint system.

\paragraph{Uniform actual coefficient comparison \(\mathrm{UC}\).}
Assume that there are \(\alpha_g>0\), \(C_{\mathrm{cmp}}\ge0\)
and \(K_0\ge1\), independent of preceding depth and of the applicable
joint structural limit, such that
\begin{equation}
 \left|\beta^{*}(D_{M,K}^{*})
                  -22\alpha_g\log K\right|
       \le C_{\mathrm{cmp}}
 \quad\hbox{for every }K=R^n\ge K_0.
 \label{lpd-UC}
\end{equation}
For the one-color convention used in the scalar comparison,
\(\alpha_g=C_H/(48\pi^2)\), with \(C_H>0\) for a nonzero probe
in the simple color factor. Only the uniform lower half of
\eqref{lpd-UC} is required for this section. Positivity in one
chosen fixed-depth limit is insufficient.

In the present assembly, Theorem~\ref{apb-main} supplies
\(\mathrm{UC}\), with \(K_0=1\), under its stated classical structural
and determinant inputs. The following argument retains the hypothesis
as a named interface so that its downstream use can be checked
independently.

\subsection{A positive leading block in every structural limit}
\label{lpd-leading-block}

Choose a positive multiple \(q\) of \(q_a\) such that
\begin{equation}
 R^q\ge K_0,\qquad
 C_{\mathrm{cmp}}\le 11\alpha_g q\log R,
 \qquad
 \delta=11\alpha_g q\log R>0.
 \label{lpd-q-choice}
\end{equation}
The exact leading cocycle and true infinite classical marginal
preservation give, in every relevant joint system,
\begin{equation}
 B_q^{*}:=\sum_{j=0}^{q-1}\beta_{j+1}^{*}(D_j^{*})
       =\beta_q^{*}(D_{0,R^q}^{*})
       \ge 22\alpha_g q\log R-C_{\mathrm{cmp}}
       \ge\delta.
 \label{lpd-infinite-block}
\end{equation}
The index zero on the right means the beginning of this block;
its preceding homogeneous depth is arbitrary. The cocycle identity
is used only in the joint infinite system. No exact additivity
of finite short-lift ratios is assumed.

Let \(\mathcal P\) be the compatible cofinal family of periods.
Every fixed power of \(R\), and every required fixed grouping,
divides all sufficiently late periods. For a \(q\)-step block
with final period \(N\), write
\begin{equation}
 M_j=MR^j,\qquad Q_j=R^{q-j}N,\qquad
 b_{0,j}^{M,N}
  =\beta_{M_jR,Q_{j+1}}(D_{M_j,R;Q_{j+1}}),\qquad
 B_q(M,N)=\sum_{j=0}^{q-1}b_{0,j}^{M,N}.
 \label{lpd-finite-block}
\end{equation}
The available joint compactness and normalized moment passage imply
that, along a subsequence of every sequence \(M_n\) and
\(N_n\in\mathcal P\) with \(N_n\to\infty\),
\begin{equation}
 B_q(M_n,N_n)\longrightarrow B_q^{*}.
 \label{lpd-finite-passage}
\end{equation}
This passage includes the long-ratio determinant in
\eqref{lpd-infinite-block}; it does not select unrelated favorable
limits. Exponential source-Hessian bounds control the second-moment
tails, and the actual finite reference denominators tend uniformly
to one.

\begin{lemma}[Uniform positive finite leading block]
\label{lpd-finite-drift}
Under \(\mathrm{UC}\), there is \(N_{\mathrm{drift}}\), independent
of the preceding depth \(M\), such that
\begin{equation}
 B_q(M,N)\ge\delta/2
 \quad\hbox{whenever }N\in\mathcal P,\quad
                         N\ge N_{\mathrm{drift}}.
 \label{lpd-finite-leading-positive}
\end{equation}
\end{lemma}
\begin{proof}
Otherwise choose violating pairs \((M_n,N_n)\) with
\(N_n\to\infty\). Joint compactness and
\eqref{lpd-finite-passage} produce a structural limit with
\(B_q^{*}\le\delta/2\), contradicting
\eqref{lpd-infinite-block}. The starting stage of a block is absorbed
into \(M\). This is a qualitative compactness argument and supplies
no numerical period threshold or convergence rate.
\end{proof}

\subsection{From leading drift to the actual finite decrement}
\label{lpd-actual-drift}

The finite leakage, zero-input parity and inherited-input comparison
give, on the established activity envelope,
\begin{align}
 |d_j(g,x_j)-b_{0,j}|
 \le{}&C_{\mathrm{int}}\omega(N_{\min})
    +C_{\mathrm{par},\beta}
          \{\tau(T(g))+g^2T(g)^6\}\notag\\
 &+C_{\mathrm{in},\beta}\,gT(g)/b .
 \label{lpd-error}
\end{align}
Here \(\omega(N)\to0\) along the compatible family, and the last
two terms tend to zero as \(g\downarrow0\), at the fixed geometric
data and field radius. The cutoff \(T(g)\) is the actual continuous
real threshold. The input factor \(b^2\) has already canceled the
\(b^{-2}\) Maxwell-extraction bound in the last term.

Choose \(N_0\) above \(N_{\mathrm{drift}}\), the reference and
finite-induction thresholds, and sufficiently large that
\begin{equation}
 C_{\mathrm{int}}\omega(N_0)\le\frac{\delta}{8q}.
 \label{lpd-period-threshold}
\end{equation}
Lower the weak-coupling ceiling \(g_{\mathrm w}>0\), if necessary,
so that for every \(0<g\le g_{\mathrm w}\),
\begin{equation}
 C_{\mathrm{par},\beta}
       \{\tau(T(g))+g^2T(g)^6\}
      +C_{\mathrm{in},\beta}gT(g)/b
       \le\frac{\delta}{8q}.
 \label{lpd-coupling-threshold}
\end{equation}
These choices are independent of the permitted history length.
The existing magnitude bound gives a fixed \(D_{\mathrm{abs}}>0\)
with \(|d_j(g,x_j)|\le D_{\mathrm{abs}}\) on this envelope.
Consequently every \(q\)-step block at arbitrary prescribed weak
couplings, started in the endpoint ball and with final period at
least \(N_0\), obeys
\begin{equation}
 B_-\le\sum_{j=k}^{k+q-1}d_j(g_j,x_j)\le B_+,\qquad
 B_-=\delta/4>0,\qquad B_+=qD_{\mathrm{abs}}.
 \label{lpd-block-bounds}
\end{equation}
Each \(q_a\)-step subblock returns to the same endpoint ball.
Thus taking \(q\) larger repeats the existing activity induction;
it does not introduce a new large-endpoint field-dilation estimate.
Neither individual leading coefficients nor individual actual
decrements need be positive.

\subsection{Finite local endpoint shooting}
\label{lpd-shooting}

\begin{theorem}[Conditional positive-drift local trajectories]
\label{lpd-main}
Assume \(\mathrm{UC}\) with the stated structural uniformity and
the established finite local inputs above. Put
\(\Gamma_{\mathrm w}=g_{\mathrm w}^{-2}\), and choose a terminal
inverse coupling
\begin{equation}
                 \Gamma_*\ge
                    \Gamma_{\mathrm w}+qD_{\mathrm{abs}}.
 \label{lpd-terminal}
\end{equation}
For every integer \(m\ge1\) and every compatible finite homogeneous
history of \(mq\) steps with final period \(N_{\mathrm{end}}\ge N_0\),
there is an initial inverse coupling
\begin{equation}
 \Gamma_*+mB_-\le\gamma_{\mathrm{bare}}
                      \le\Gamma_*+mB_+
 \label{lpd-initial-interval}
\end{equation}
whose actual local extraction trajectory reaches \(\Gamma_*\)
and stays in the weak range throughout. Its initial period is
\(R^{mq}N_{\mathrm{end}}\); all intermediate periods obey the
same compatibility requirements. Its activity classes obey the
same endpoint and intermediate \(O(b^2)\) bounds.
\end{theorem}

\begin{proof}
For the proof alone, evaluate the local map at
\begin{equation}
 g_{\mathrm{eval}}(\gamma)
             =\max\{\gamma,\Gamma_{\mathrm w}\}^{-1/2},
 \label{lpd-clipped-coupling}
\end{equation}
while retaining the unclipped scalar update
\(\gamma_{\mathrm{next}}
 =\gamma-d_j(g_{\mathrm{eval}}(\gamma),x)\).
The evaluated couplings are always weak, so finite induction and
\eqref{lpd-block-bounds} apply to every auxiliary history, regardless
of the scalar value. The prescribed-coupling estimates apply to
its realized sequence; no adaptation changes their constants.

For this fixed finite regulator the joint coupling/class continuity
of the normalized local map makes the terminal scalar continuous
in its initial value. On the compact interval
\eqref{lpd-initial-interval}, the absolute increment bound keeps
all intermediate scalars bounded. Their evaluated couplings
therefore lie in a compact positive interval. No continuity at
\(g=0\), volume-uniform continuity modulus, or continuous choice
of equivalent activity representatives is needed.

At the left endpoint of \eqref{lpd-initial-interval}, the terminal
scalar is at most \(\Gamma_*\), by the lower block bound. At its
right endpoint it is at least \(\Gamma_*\), by the upper block
bound. The intermediate value theorem selects a clipped history
ending exactly at \(\Gamma_*\). Summing its future block decrements
backwards gives
\begin{equation}
 \Gamma_*+(m-k)B_-\le\gamma_{kq}
                   \le\Gamma_*+(m-k)B_+,
 \qquad 0\le k\le m.
 \label{lpd-block-endpoints}
\end{equation}
Inside each block,
\begin{equation}
 \gamma_{kq+r}\ge\gamma_{kq}-rD_{\mathrm{abs}}
       \ge\Gamma_*-qD_{\mathrm{abs}}
       \ge\Gamma_{\mathrm w},
 \qquad 0\le r\le q.
 \label{lpd-clipping-absent}
\end{equation}
Clipping is therefore absent from the selected history.
Its actual updates are exactly \eqref{lpd-updates} with
\(g_j=\gamma_j^{-1/2}\). The finite quotient relocalizations
replace representatives only by equal normalized totals; their
scalar differences are carried by the exact vacuum factors.
Thus the conclusion concerns the genuine local density update,
not an approximate scalar recurrence.
\end{proof}

The bare coupling in this theorem satisfies
\begin{equation}
 (\Gamma_*+mqD_{\mathrm{abs}})^{-1/2}
       \le g_{\mathrm{bare}}
       \le(\Gamma_*+m\delta/4)^{-1/2}.
 \label{lpd-bare-asymptotic}
\end{equation}
It is of order \(m^{-1/2}\), equivalently of order
\((mq)^{-1/2}\) with \(q\) fixed, and reaches the same positive
weak endpoint \(g_*=\Gamma_*^{-1/2}\). The theorem asserts neither
uniqueness of the tuned initial value nor consistency of the
choices at different horizons.

\subsection{Parameter order and scope}
\label{lpd-scope}

First fix \(R,J,b\) and the available \(q_a\)-block activity system.
The uniform coefficient comparison determines a sufficiently large
fixed multiple \(q\) by \eqref{lpd-q-choice}. Then choose the
period threshold \(N_0\), the weak ceiling \(g_{\mathrm w}\), and
the fixed terminal margin \(\Gamma_*\). The horizon \(m\) is
arbitrary after these choices. No quantitative finite-volume rate
and no exact finite classical marginal identity enter this order.

The only new coefficient input of this section is
\eqref{lpd-UC}. If an actual comparison theorem supplies it in
every required joint structural system, this section discharges
the positive-drift hypothesis of finite local shooting. The
comparison is supplied here by Theorem~\ref{apb-main}. The
conclusion remains about homogeneous local analytic histories
and their justified continuation. It does not identify those
histories with the full Wilson measure, control mixed or large-field
sectors, construct continuum observable measures, prove a
nontrivial continuum Yang--Mills theory, or establish a mass gap.
